\begin{document}
\setlength{\parindent}{0pt}
\setlength{\parskip}{4pt}

% ---- Cover: the unit's cover art, full bleed, in colour ------------------
% The art is 2:3 and A4 is taller-than-wide by less, so it fills the width and
% loses a little off the bottom; the baked-in title at the top stays whole.
\pdfbookmark[1]{Cover}{cover}
\thispagestyle{empty}
\stcover
\begin{tikzpicture}[remember picture, overlay]
  
  % Name box, top right, clear of the title.
  \node[anchor=north east, fill=white, draw=Ink, line width=0.6pt, rounded corners=1.5mm,
        inner xsep=2.6mm, inner ysep=2.2mm, text width=40mm, font=\headingfont\small]
    at ([xshift=-6mm, yshift=-6mm]current page.north east)
    {\textbf{Name}\par\vspace{5mm}\hrule\vspace{2.6mm}\textbf{Class}\par\vspace{5mm}\hrule};
  % What this book is: the cover art says Workbook, so name the booklet here.
  \node[anchor=south east, fill=white, draw=Ink, line width=0.6pt, rounded corners=1.5mm,
        inner xsep=3.2mm, inner ysep=2.4mm, font=\headingfont\bfseries\large, align=right]
    at ([xshift=-6mm, yshift=8mm]current page.south east)
    {Experiment Booklet\\[-1pt]{\normalfont\headingfont\small Experiments 1--9}};
\end{tikzpicture}
\mbox{}
\clearpage

% ---- Inside cover ---------------------------------------------------------
\thispagestyle{plain}
{\headingfont\bfseries\fontsize{24}{28}\selectfont Separation Technologies and Techniques}\par\vspace{1mm}
{\headingfont\Large Experiment Booklet \textbullet\ Year 7}\par\vspace{1mm}
{\small\itshape VC 2.0: VC2S8U05 \textbullet\ VC2S8U06 \textbullet\ VC2S8U07 \textbullet\ VC2S8I01--I08 \textbullet\ VC2S8H02}\par\vspace{6mm}
\renewcommand{\arraystretch}{2.2}
\begin{tabular}{@{}l p{130mm}@{}}
\textbf{Name} & \hrulefill \\ \textbf{Class} & \hrulefill \\ \textbf{Teacher} & \hrulefill \\ \textbf{Lab group} & \hrulefill \\
\end{tabular}
\renewcommand{\arraystretch}{1}

\begin{countrybox}{Acknowledgement of Country}
This booklet was created on the unceded Country of the Wurundjeri Woi-wurrung people of the Kulin Nation. We acknowledge Elders past, present and emerging. Aboriginal and Torres Strait Islander peoples have observed, tested and passed on knowledge of Country for tens of thousands of years: science in the way this booklet teaches it, by careful observation, measurement and sharing what was found.
\end{countrybox}

\begin{infobox}{How this booklet works}
This booklet holds every experiment in the unit. Bring it to every practical lesson. Each experiment opens with its key vocabulary, learning intentions and skills focus, then runs \emph{Technique} $\rightarrow$ \emph{Plan and safety} $\rightarrow$ \emph{Method} $\rightarrow$ \emph{Results} $\rightarrow$ \emph{Analysis and evaluation} $\rightarrow$ \emph{Skill check}. Each kind of box has its own frame and mark:\par\smallskip
\renewcommand{\arraystretch}{1.35}
\begin{tabular}{@{}p{32mm}l@{}}
\textbf{Technique} \tagicon{\icokey} & How to do it properly (thin frame) \\
\textbf{I do} \tagicon{\icodots{1}} & Your teacher works an example through --- you follow (solid frame) \\
\textbf{We do} \tagicon{\icodots{2}} & You try one together, with help (dashed frame) \\
\textbf{You do} \tagicon{\icodots{3}} & On your own: check your understanding (heavy frame) \\
\textbf{Going further} \tagicon{\icoup} & Above level: Year 8 extension (double frame) \\
\textbf{Skills focus} \tagicon{\icoskill} & The science skills the experiment builds \\
\end{tabular}

\smallskip
Every skill comes from the faculty \textbf{Science Skills} list (the toolkit on page \pageref{sec:toolkit}). A skill always carries the same code, for example \sk{PC4} or \sk{PA1}, so you can follow it from this unit into every other unit this year.
\end{infobox}

\clearpage
{\headingfont\bfseries\LARGE Contents}\par\vspace{1mm}\noindent\rule{\linewidth}{0.6pt}\par\vspace{2mm}
\begin{ctab}
\renewcommand{\arraystretch}{1.5}
\begin{tabular}{|L{2.8cm}|L{10.2cm}|C{1.6cm}|}
  \hline\thead{Section} & \thead{Title} & \thead{Page} \\ \hline
  Skills & \hyperlink{toolkit}{The Science Skills toolkit} & \hyperlink{toolkit}{\pageref{sec:toolkit}} \\ \hline
  Skills & \hyperlink{tracker}{My skills tracker} & \hyperlink{tracker}{\pageref{sec:tracker}} \\ \hline
  Unit & \hyperlink{glance}{Separation techniques at a glance} & \hyperlink{glance}{\pageref{sec:glance}} \\ \hline
  Experiment 1 & \hyperlink{exp.1}{Solubility (Lesson 3)} & \hyperlink{exp.1}{\pageref{sec:1}} \\ \hline
  Experiment 2 & \hyperlink{exp.2}{Surface Area and Dissolving (Lesson 6)} & \hyperlink{exp.2}{\pageref{sec:2}} \\ \hline
  Experiment 3 & \hyperlink{exp.3}{Fast and Slow Evaporation (Lesson 9)} & \hyperlink{exp.3}{\pageref{sec:3}} \\ \hline
  Experiment 4 & \hyperlink{exp.4}{Rates of Evaporation (Lesson 11)} & \hyperlink{exp.4}{\pageref{sec:4}} \\ \hline
  Experiment 5 & \hyperlink{exp.5}{Solubility and Temperature (Lesson 13)} & \hyperlink{exp.5}{\pageref{sec:5}} \\ \hline
  Experiment 6 & \hyperlink{exp.6}{Growing Copper Sulfate Crystals (Lesson 14)} & \hyperlink{exp.6}{\pageref{sec:6}} \\ \hline
  Experiment 7 & \hyperlink{exp.7}{Filtration (Lesson 15)} & \hyperlink{exp.7}{\pageref{sec:7}} \\ \hline
  Experiment 8 & \hyperlink{exp.8}{Distillation (Lesson 17)} & \hyperlink{exp.8}{\pageref{sec:8}} \\ \hline
  Experiment 9 & \hyperlink{exp.9}{The Great Green Mix Up (Lesson 18)} & \hyperlink{exp.9}{\pageref{sec:9}} \\ \hline
  Glossary & \hyperlink{glossary}{Key vocabulary, with the page each word first appears on} & \hyperlink{glossary}{\pageref{sec:glossary}} \\ \hline
\end{tabular}
\end{ctab}

% ══════════════════════════════════════════════════════════════════════════════
%  THE SCIENCE SKILLS TOOLKIT — faculty continuum, Years 7–8 column
% ══════════════════════════════════════════════════════════════════════════════
\clearpage
\hypertarget{toolkit}{}\pdfbookmark[1]{The Science Skills toolkit}{toolkitbm}
\markboth{The Science Skills toolkit}{}{\headingfont\bfseries\LARGE The Science Skills toolkit}\label{sec:toolkit}\par\vspace{1mm}\noindent\rule{\linewidth}{0.6pt}\par
These are the skills scientists use in every investigation. They come from the Science Inquiry strand of the Victorian Curriculum, grouped the way our faculty teaches them. Each experiment in this booklet tells you which skills it builds.

\newcommand{\tkhead}[2]{\multicolumn{3}{|l|}{\cellcolor{black!10}\rule{0pt}{4.6mm}\headingfont\bfseries #2} \\ \hline}
\newcommand{\tkrow}[1]{\skc{#1}\enspace{\small\bfseries\skname{#1}} & \skdesc{#1} & {\small\skcode{#1}} \\ \hline}
\renewcommand{\arraystretch}{1.35}
\begin{ctab}
\begin{longtable}{|W{5.0cm}|W{9.0cm}|Z{1.7cm}|}
  \hline\thead{Skill} & \thead{What it means} & \thead{Code} \\ \hline
  \endhead
  \tkhead{Teal}{Questioning and predicting}
  \tkrow{QP1}
  \tkrow{QP2}
  \tkhead{Blue}{Planning and conducting}
  \tkrow{PC1}
  \tkrow{PC2}
  \tkrow{PC3}
  \tkrow{PC4}
  \tkrow{PC5}
  \tkrow{PC6}
  \tkrow{PC7}
  \tkrow{PC8}
  \tkhead{GoldDark}{Processing, modelling and analysing}
  \tkrow{PA1}
  \tkrow{PA2}
  \tkrow{PA3}
  \tkrow{PA4}
  \tkrow{PA5}
  \tkrow{PA6}
  \tkrow{PA7}
  \tkrow{PA8}
  \tkhead{OrangeDark}{Evaluating}
  \tkrow{EV1}
  \tkrow{EV2}
  \tkrow{EV3}
  \tkrow{EV4}
  \tkrow{EV5}
  \tkhead{PurpleDark}{Communicating}
  \tkrow{CO1}
  \tkrow{CO2}
\end{longtable}
\end{ctab}
\renewcommand{\arraystretch}{1}

\begin{infobox}{Writing like a scientist \textbar\ CO1 \skname{CO1}}
Scientists write in the \textbf{third person}: say what was done, not who did it. Write \emph{``The water was heated for 10 minutes''}, not \emph{``I heated the water for 10 minutes''}. Keep sentences short, spell scientific words correctly, and always give a unit with a measurement.
\end{infobox}

% ══════════════════════════════════════════════════════════════════════════════
%  SKILLS TRACKER
% ══════════════════════════════════════════════════════════════════════════════
\clearpage
\hypertarget{tracker}{}\pdfbookmark[1]{My skills tracker}{trackerbm}
\markboth{My skills tracker}{}{\headingfont\bfseries\LARGE My skills tracker}\label{sec:tracker}\par\vspace{1mm}\noindent\rule{\linewidth}{0.6pt}\par
A dot \textbullet{} shows the experiments that build each skill. After each experiment, write how confident you are with each of its skills in its column: \textbf{1}~not yet, \textbf{2}~nearly, \textbf{3}~got it.

\newcommand{\dotc}{{\large\textbullet}}
\newcommand{\trow}[2]{\rule[-0.3cm]{0pt}{0.82cm}\skc{#1}\enspace{\footnotesize\skname{#1}} & #2 \\ \hline}
\setlength{\tabcolsep}{3pt}
\begin{ctab}
\begin{tabular}{|W{6.1cm}|*{9}{Z{0.95cm}|}}
  \hline\thead{Skill} & \thead{1} & \thead{2} & \thead{3} & \thead{4} & \thead{5} & \thead{6} & \thead{7} & \thead{8} & \thead{9} \\ \hline
  \trow{QP1}{\dotc & \dotc &  & \dotc &  &  &  &  & }
  \trow{QP2}{ &  &  & \dotc &  &  &  &  & \dotc}
  \trow{PC1}{ &  &  & \dotc &  &  &  &  & \dotc}
  \trow{PC2}{\dotc & \dotc &  & \dotc &  &  & \dotc &  & }
  \trow{PC3}{ &  & \dotc &  &  & \dotc &  & \dotc & \dotc}
  \trow{PC4}{ &  &  & \dotc & \dotc &  &  &  & }
  \trow{PC5}{ &  &  &  &  & \dotc &  &  & }
  \trow{PC6}{\dotc &  & \dotc &  & \dotc & \dotc & \dotc & \dotc & }
  \trow{PC7}{ &  &  &  &  &  & \dotc & \dotc & \dotc}
  \trow{PC8}{ & \dotc &  &  &  &  &  &  & }
  \trow{PA1}{\dotc & \dotc &  &  & \dotc &  &  &  & }
  \trow{PA2}{ &  &  & \dotc & \dotc &  &  &  & }
  \trow{PA3}{ & \dotc &  & \dotc &  &  &  &  & }
  \trow{PA5}{ & \dotc &  &  & \dotc &  &  &  & }
  \trow{PA6}{ &  & \dotc &  &  &  &  &  & }
  \trow{PA7}{ &  &  &  & \dotc &  &  &  & }
  \trow{EV1}{ & \dotc &  &  &  &  &  &  & }
  \trow{EV2}{ &  &  &  & \dotc &  &  & \dotc & }
  \trow{EV3}{ & \dotc &  & \dotc &  &  & \dotc &  & \dotc}
  \trow{EV5}{ &  & \dotc &  &  & \dotc &  &  & }
  \trow{CO1}{\dotc &  &  &  &  &  &  &  & \dotc}
  \trow{CO2}{ &  & \dotc &  &  & \dotc & \dotc & \dotc & }
\end{tabular}
\end{ctab}
\setlength{\tabcolsep}{6pt}
{\small\color{DarkGrey} Columns 1--9 are the experiments. The other skills are built in the experiment booklets for other units.}

% ══════════════════════════════════════════════════════════════════════════════
%  SEPARATION TECHNIQUES AT A GLANCE
% ══════════════════════════════════════════════════════════════════════════════
\clearpage
\hypertarget{glance}{}\pdfbookmark[1]{Separation techniques at a glance}{glancebm}
\markboth{Separation techniques at a glance}{}{\headingfont\bfseries\LARGE Separation techniques at a glance}\label{sec:glance}\par\vspace{1mm}\noindent\rule{\linewidth}{0.6pt}\par
Every separation uses a property that one part of a mixture has and the others do not. Use this page to plan a separation, and tick each technique when you have used it.
\renewcommand{\arraystretch}{1.45}
\begin{ctab}
\begin{tabular}{|L{2.9cm}|L{3.6cm}|L{4.4cm}|L{2.5cm}|C{1.3cm}|}
  \hline\thead{Technique} & \thead{Property it uses} & \thead{Separates} & \thead{Experiment} & \thead{Used} \\ \hline
  Sieving & Size of the pieces & Large solids from small solids & 9 & \tickbox \\ \hline
  Magnetic separation & Magnetic or not & Iron and steel from other solids & 9 & \tickbox \\ \hline
  Decanting & Insoluble solid sinks & A liquid from a settled solid & 7 & \tickbox \\ \hline
  Filtration & Insoluble; larger than the pores & An insoluble solid from a liquid & 7, 9 & \tickbox \\ \hline
  Evaporation & Solute does not evaporate & A dissolved solid from its solvent & 3, 4, 9 & \tickbox \\ \hline
  Crystallisation & Solubility falls as solvent leaves & Pure crystals from a solution & 3, 6 & \tickbox \\ \hline
  Distillation & Different boiling points & A pure liquid from a solution & 8 & \tickbox \\ \hline
  Chromatography & Different solubility in the solvent & The dyes in an ink & Lesson & \tickbox \\ \hline
\end{tabular}
\end{ctab}
\renewcommand{\arraystretch}{1}

\begin{infobox}{Pure substance or mixture?}
A \textbf{pure substance} contains only one kind of particle: an element (such as copper) or a compound (such as water). A \textbf{mixture} contains two or more substances that are not chemically joined, so it can be separated by physical methods like the ones above. A \textbf{solution} is a mixture in which a solute is dissolved in a solvent.
\end{infobox}

